\VerdierPhysicalPageAnchor{0046}%
\VerdierSemanticAnchor{ch1-1-4-3-canonical}
This operation is called the \emph{canonical operation}. We propose to
study the $H$-functors

\VerdierSemanticAnchor{ch1-1-4-3-eq3}
\begin{equation}
  (F,m)\colon(\mathcal C,H,T_f,c_f)
  \longrightarrow(\mathcal C',H,T'_f,c'_f)
  \tag{1.4.3.3}
\end{equation}
that induce the canonical operation of the group
$\mathbb Z/2\mathbb Z$ on the functor $F$.

\VerdierSemanticAnchor{ch1-1-4-3-extension}
Let $\dot\alpha\in H^2(G,\mathbb Z/2\mathbb Z)$ be the element
describing the extension

\VerdierSemanticAnchor{ch1-1-4-3-eq4}
\begin{equation}
  1\longrightarrow\mathbb Z/2\mathbb Z
  \xrightarrow{n}H\xrightarrow{f}G\longrightarrow e.
  \tag{1.4.3.4}
\end{equation}

\VerdierSemanticAnchor{ch1-1-4-3-cocycle}
Let $\alpha_s\in Z^2(G,\mathbb Z/2\mathbb Z)$ be a normalized
inhomogeneous cocycle in the class $\dot\alpha$, defined by a normalized
section $s$ of $f$. We have

\VerdierSemanticAnchor{ch1-1-4-3-eq5}
\begin{equation}
  \begin{aligned}
    \alpha_s(\sigma_1,\sigma_2)
      &=s(\sigma_2)s(\sigma_1\sigma_2)^{-1}s(\sigma_1),
      & \sigma_1,\sigma_2&\in G,\\
    s(e)&=\text{the identity element of }H,\\
    \alpha_s(\sigma_1,e)&=\alpha_s(e,\sigma_1)=1,\\
    \delta\alpha_s(\sigma_1,\sigma_2,\sigma_3)
      &=\alpha_s(\sigma_2,\sigma_3)
        \alpha_s(\sigma_1\sigma_2,\sigma_3)^{-1}\\
      &\quad{}\alpha_s(\sigma_1,\sigma_2\sigma_3)
        \alpha_s(\sigma_1,\sigma_2)^{-1}=1,
      & \sigma_1,\sigma_2,\sigma_3&\in G.
  \end{aligned}
  \tag{1.4.3.5}
\end{equation}

\VerdierSemanticAnchor{ch1-1-4-3-bijection}
The section $s\colon G\to H$ defines an isomorphism of sets
\[
  H\xrightarrow{\sim}\mathbb Z/2\mathbb Z\times G,
\]
and by transporting the group law of $H$ to
$\mathbb Z/2\mathbb Z\times G$ through this isomorphism, we obtain the
group law

\VerdierSemanticAnchor{ch1-1-4-3-eq6}
\begin{equation}
  (\varepsilon_1,\sigma_1)(\varepsilon_2,\sigma_2)
  =\bigl(\varepsilon_1\varepsilon_2
    \alpha_s(\sigma_1,\sigma_2),\sigma_1\sigma_2\bigr).
  \tag{1.4.3.6}
\end{equation}

\VerdierSemanticAnchor{ch1-1-4-3-identification}
Let us now identify $H$ with the group described by (1.4.3.6). For
every $(\varepsilon_1,\sigma_1),(\varepsilon_2,\sigma_2)\in H$, we have

\VerdierSemanticAnchor{ch1-1-4-3-eq7}
\begin{equation}
  \begin{aligned}
    T_f(\varepsilon_1,\sigma_1)&=T(\sigma_1),\\
    T'_f(\varepsilon_1,\sigma_1)&=T'(\sigma_1),\\
    c_f\bigl((\varepsilon_1,\sigma_1),(\varepsilon_2,\sigma_2)\bigr)
      &=c(\sigma_1,\sigma_2),\\
    c'_f\bigl((\varepsilon_1,\sigma_1),(\varepsilon_2,\sigma_2)\bigr)
      &=c'(\sigma_1,\sigma_2)
  \end{aligned}
  \tag{1.4.3.7}
\end{equation}
